\chapter{Bases de Datos Relacionales, Integridad Formal y Lenguaje DDL}
\label{cap:modelo_relacional_ddl}

\section{Marco Conceptual y Te\'orico (Curr\'iculo Universitario)}

\subsection{Fundamentos del Modelo Relacional}
Propuesto por Edgar F. Codd (1970), el modelo se sustenta en la teor\'ia matem\'atica de conjuntos y la l\'ogica de predicados de primer orden:
\begin{itemize}
    \item \textbf{Relaci\'on:} Una tabla bidimensional de datos donde el orden de filas y columnas es irrelevante.
    \item \textbf{Tupla:} Una fila individual en la relaci\'on que representa una instancia \'unica de la entidad.
    \item \textbf{Atributo:} Columna con nombre y tipo de dato formal.
    \item \textbf{Dominio:} Conjunto finito o infinito de valores at\'omicos permitidos para un atributo.
\end{itemize}

\subsection{Restricciones de Integridad Formales}
\begin{enumerate}
    \item \textbf{Integridad de Entidad:} Toda tabla debe poseer una Clave Primaria (\textit{Primary Key}) \'unica y no nula (\texttt{NOT NULL}).
    \item \textbf{Integridad Referencial:} Toda Clave For\'anea (\textit{Foreign Key}) debe coincidir con un valor existente de clave primaria en la tabla referenciada o ser nula.
    \item \textbf{Integridad de Dominio:} Los valores deben adherirse a las restricciones de tipo y chequeos (\texttt{CHECK}).
\end{enumerate}

---

\section{Casos de Estudio y Pr\'actica Aplicada}

\begin{lstlisting}[style=sqlstyle, caption={Definici\'on DDL con restricciones de integridad completas}]
-- Creaci\'on de esquema robusto para salas de cine y funciones
CREATE TABLE cinemas (
    cinema_id SERIAL PRIMARY KEY,
    cinema_name VARCHAR(100) NOT NULL,
    city VARCHAR(50) NOT NULL,
    total_screens SMALLINT DEFAULT 1 CHECK (total_screens > 0),
    is_active BOOLEAN DEFAULT TRUE,
    created_at TIMESTAMPTZ DEFAULT CURRENT_TIMESTAMP
);

CREATE TABLE movie_showtimes (
    showtime_id BIGSERIAL PRIMARY KEY,
    cinema_id INTEGER NOT NULL,
    movie_title VARCHAR(150) NOT NULL,
    start_time TIMESTAMPTZ NOT NULL,
    ticket_price NUMERIC(8, 2) NOT NULL CHECK (ticket_price >= 0.00),
    CONSTRAINT fk_showtimes_cinema
        FOREIGN KEY (cinema_id)
        REFERENCES cinemas (cinema_id)
        ON DELETE CASCADE
        ON UPDATE CASCADE
);

-- \'indice para optimizaci\'on de b\'usqueda por horario y cine
CREATE INDEX idx_showtimes_cinema_time ON movie_showtimes (cinema_id, start_time);
\end{lstlisting}

---

\section{Taller de Consolidaci\'on del Cap\'itulo}

\begin{businessproblem}
Implementar un m\'odulo DDL para una tienda de videojuegos en l\'inea, asegurando que los descuentos no superen el 100\% del valor unitario del producto y que las categor\'ias mantengan nombres \'unicos en min\'usculas.
\end{businessproblem}

\subsection{Soluci\'on T\'ecnica en PostgreSQL / MySQL}

\begin{lstlisting}[style=sqlstyle, caption={Script DDL con constraints y tipos personalizados}]
CREATE TABLE game_categories (
    category_id SERIAL PRIMARY KEY,
    category_slug VARCHAR(50) NOT NULL UNIQUE,
    display_name VARCHAR(80) NOT NULL,
    CONSTRAINT chk_category_slug_lowercase CHECK (category_slug = LOWER(category_slug))
);

CREATE TABLE video_games (
    game_id BIGSERIAL PRIMARY KEY,
    category_id INTEGER NOT NULL REFERENCES game_categories(category_id),
    title VARCHAR(120) NOT NULL,
    base_price NUMERIC(10, 2) NOT NULL CHECK (base_price >= 0),
    discount_percentage NUMERIC(5, 2) DEFAULT 0.00 CHECK (discount_percentage BETWEEN 0.00 AND 100.00),
    release_date DATE NOT NULL,
    is_available BOOLEAN DEFAULT TRUE
);
\end{lstlisting}

\begin{engtip}
\begin{itemize}
    \item \textbf{Restricciones de Dominio con \texttt{CHECK}:} Garantizan que datos corruptos o valores negativos jam\'as ingresen al motor de almacenamiento, independientemente del lenguaje o framework utilizado en el backend.
    \item \textbf{\texttt{ON DELETE CASCADE}:} Debe usarse con precauci\'on extrema en entornos de producci\'on financiera; en esquemas transaccionales se prefiere \texttt{ON DELETE RESTRICT} o borrado l\'ogico mediante banderas (\textit{soft delete}).
\end{itemize}
\end{engtip}
